\documentclass[11pt,a4paper]{article}
\usepackage[T1]{fontenc}
\usepackage[utf8]{inputenc}
\usepackage{lmodern,amsmath,amssymb}
\usepackage[margin=25mm]{geometry}
\usepackage{longtable,booktabs,array,calc}
\usepackage[hidelinks]{hyperref}
\hypersetup{pdftitle={Flowing before decimating leaves every norm-based criterion unchanged: what is needed is a},pdfauthor={navstokgap research working notes}}
\newcommand{\tightlist}{\setlength{\itemsep}{0pt}\setlength{\parskip}{0pt}}
\setlength{\emergencystretch}{3em}
\setcounter{secnumdepth}{0}
\begin{document}
\hypertarget{flowing-before-decimating-leaves-every-norm-based-criterion-unchanged-what-is-needed-is-a-statement-about-the-state}{%
\section{Flowing before decimating leaves every norm-based criterion
unchanged: what is needed is a statement about the
state}\label{flowing-before-decimating-leaves-every-norm-based-criterion-unchanged-what-is-needed-is-a-statement-about-the-state}}

Conjugating by the flow before blocking does not improve the criterion
of \href{blocking-criterion-monotone.md}{the blocking-criterion note},
and the reason is one line: \(\Phi_t\) is a bijection of the
configuration space, so
\[\big\|V\circ\Phi_t\big\|_\infty=\sup_U V\big(\Phi_t(U)\big)
=\sup_{U'}V(U')=\big\|V\big\|_\infty,\] even though
\(V(\Phi_t(U))\le V(U)\) holds pointwise by the monotonicity of the
flow. The straddling-plaquette norm that enters \(\beta_{\rm block}\) is
therefore unchanged, the block gap is unchanged because conjugation is
unitary, and the factor \(3M^2/8\) by which blocking loses is unchanged.
The same argument disposes of every criterion built from operator norms:
\textbf{no reordering of conjugation and decimation can help, because
the quantities such a criterion uses are conjugation-invariant or
bijection-invariant.} The flow's gain is pointwise and therefore visible
only in expectation, that is, only against a state whose weight sits
where the flowed action is small. With that, the last inexpensive idea
in this programme is closed, and the residual statement is sharp: the
remaining input is control of the ground-state measure, which is the
constructive problem in its Euclidean form. Constants explicit; nothing
promoted.

\hypertarget{the-three-invariances}{%
\subsection{1. The three invariances}\label{the-three-invariances}}

Let \(\Phi_t\) be the lattice flow map and \(W_t\) its unitary
implementation, as in \href{flow-conjugation-truncation.md}{the
flow-conjugation note} Proposition 1.

\textbf{(i) Spectra are conjugation-invariant.} Every eigenvalue of
\(H\), every gap, and in particular the block gap \(\Delta_M\) of a
block Hamiltonian, is unchanged by \(H\mapsto W_tHW_t^*\).

\textbf{(ii) Sup norms of multiplication operators are
bijection-invariant.} For \(f\in C(\mathcal M)\) and a bijection
\(\Phi\) of \(\mathcal M\), \(\|f\circ\Phi\|_\infty=\|f\|_\infty\).
Applied to the magnetic term,
\[\big\|V\circ\Phi_t\big\|_\infty=\big\|V\big\|_\infty
=\frac{4N\hbar c}{ag^2}\cdot\#\{\text{plaquettes}\},\] and to the
straddling part, \(\|\sum_{p\ \rm str}w_p\circ\Phi_t\|_\infty\) is
bounded by the same \(6M^2\cdot4N\hbar c/(ag^2)\) per block as before.
The pointwise inequality \(V\circ\Phi_t\le V\), which the monotonicity
of the flow supplies, is compatible with equality of the suprema because
\(\Phi_t\) maps the configuration space onto itself.

\textbf{(iii) Hence the criterion is unchanged.} The blocking criterion
\[\beta_{\rm block}=\frac{\big\|\sum_{p\ \rm str}w_p\big\|_\infty}{\Delta_M}
=\frac{24M^2N}{g^2\,\delta(g;M)}\] has a numerator fixed by (ii) and a
denominator fixed by (i), so conjugating before blocking leaves it
exactly as computed in \href{blocking-criterion-monotone.md}{the
blocking-criterion note}, and blocking still loses by \(3M^2/8\).

\hypertarget{why-no-reordering-helps}{%
\subsection{2. Why no reordering helps}\label{why-no-reordering-helps}}

Any criterion of the form (perturbation norm)/(unperturbed gap) is built
from two quantities of which the first is invariant under bijections of
the configuration space and the second under unitary conjugation. The
flow supplies exactly a bijection and a unitary. So:

\textbf{Proposition.} Let \(\mathcal C(H_0,\phi)\) be any criterion
depending on \(H_0\) only through its spectrum and on \(\phi\) only
through \(\|\phi\|_\infty\). Then \(\mathcal C\) takes the same value
for \((H_0,\phi)\) and for \((W_tH_0W_t^*,\,\phi\circ\Phi_t)\), for
every flow time \(t\).

This covers the Yarotsky criterion used in
\href{strong-coupling-uniform-gap.md}{the T2 note}, its blocked version,
and the Schur-complement criteria of
\href{weak-coupling-feshbach-reduction.md}{the Feshbach note}, all of
which reduce to comparisons of norms with gaps.

\hypertarget{where-the-flows-gain-actually-lives}{%
\subsection{3. Where the flow's gain actually
lives}\label{where-the-flows-gain-actually-lives}}

The flow does reduce the action: \(V(\Phi_t(U))\le V(U)\) for every
\(U\), with strict inequality away from critical points, and
\(\frac{d}{ds}S(B_s)=-\|D^*G\|_2^2\). That gain appears in expectations,
\[\big\langle\psi,\;(V\circ\Phi_t)\,\psi\big\rangle\ \le\ \big\langle\psi,\;V\psi\big\rangle
\qquad\text{for every state }\psi,\] and is large when the weight of
\(\psi\) sits where the flow moves the configuration far, that is, on
rough configurations. So the flow converts a rough state into a smooth
one at the level of expectations while leaving every worst-case bound
alone.

A criterion that could see the gain must therefore be \textbf{relative
to the state}, of the form
\[\big\langle\psi,\phi\,\psi\big\rangle\ \le\ \epsilon\,\big\langle\psi,(H-E_0)\psi\big\rangle+\eta\|\psi\|^2
\qquad\text{on the low-energy subspace},\] which is the relative form
already isolated in \href{schur-error-ultraviolet.md}{the Schur note}
Lemma 1\('\) and in the relative hypothesis of
\href{weak-coupling-feshbach-reduction.md}{the Feshbach note} §3a. The
quantity it needs is the distribution of the field strength in the
low-energy states, that is, the ground-state measure
\(|\Omega(U)|^2\,dU\), which is a probability measure on the
configuration space and is exactly the object the Euclidean formulation
supplies explicitly as \(e^{-S}dU\) and the Hamiltonian formulation
leaves implicit.

\hypertarget{the-residual-statement}{%
\subsection{4. The residual statement}\label{the-residual-statement}}

Collecting the closed routes of this programme:

\begin{longtable}[]{@{}
  >{\raggedright\arraybackslash}p{(\columnwidth - 4\tabcolsep) * \real{0.3333}}
  >{\raggedright\arraybackslash}p{(\columnwidth - 4\tabcolsep) * \real{0.3333}}
  >{\raggedright\arraybackslash}p{(\columnwidth - 4\tabcolsep) * \real{0.3333}}@{}}
\toprule\noalign{}
\begin{minipage}[b]{\linewidth}\raggedright
method
\end{minipage} & \begin{minipage}[b]{\linewidth}\raggedright
status
\end{minipage} & \begin{minipage}[b]{\linewidth}\raggedright
reason
\end{minipage} \\
\midrule\noalign{}
\endhead
\bottomrule\noalign{}
\endlastfoot
variational upper bounds & closed for \(m>0\) & a spectral measure with
all negative moments finite can have \(m=0\) \\
expansion around the free theory & closed & bound reaches \(m\) only at
the confinement scale \\
blocking by projection & closed & boundary grows like \(M^2\), block gap
does not \\
conjugation by the flow & closed & spectrum-invariant; the gain is the
truncation, which is cheap \\
flow before decimation & closed (this note) & norms are
bijection-invariant \\
decimation with controlled couplings & \textbf{open} & the constructive
problem \\
\end{longtable}

\textbf{Addendum.} The ground-state measure estimate asked for in
Section 5 exists: \href{agmon-ground-state-suppression.md}{the Agmon
note} proves the Agmon identity for the Kogut--Susskind ground state and
gives a suppression \(e^{-\lambda n}\) of configurations with \(n\)
excess plaquettes, with \(\lambda\simeq2\sqrt{2v}/g^2\), the coupling
dependence of the Euclidean Wilson weight.

Every method that uses only norms and spectra is exhausted, and each was
closed by a computation with explicit constants. What is left needs the
distribution of the field in the low-energy states, and that is where
the Euclidean construction has its tools and the Hamiltonian formulation
has none of its own.

\hypertarget{consequence-for-state}{%
\subsection{5. Consequence for STATE}\label{consequence-for-state}}

The question raised at the end of
\href{lattice-truncation-uniform.md}{the lattice-truncation note} is
answered in the negative, with the invariance that makes it so. The
programme's norm-based line is complete. The next thing that would
advance it is an estimate on the ground-state measure of the
Kogut--Susskind Hamiltonian at intermediate coupling, for instance a
bound of the form \(|\Omega(U)|^2\le C\,e^{-\lambda S_w(U)}\) with
\(\lambda>0\) uniform in the volume, which would import the Euclidean
large-field machinery into the Hamiltonian setting. That is a well-posed
question and it is the natural successor to everything above.
\end{document}
